% ============================================================
% Code-Review Checklist — Tutorial 3 (code walkthrough)
% Compile: tectonic -X compile code_review_checklist.tex
% ============================================================
\documentclass[11pt]{article}
\usepackage[margin=2.3cm]{geometry}
\usepackage{amsmath,amssymb}
\usepackage{booktabs}
\usepackage{fontspec}
\usepackage{microtype}
\usepackage{xcolor}
\usepackage[hidelinks]{hyperref}
\usepackage{enumitem}

\setmainfont{Times New Roman}
\newcommand{\e}[1]{{\emojifont #1}}

\setlength{\parindent}{0pt}
\setlength{\parskip}{0.5em}
\renewcommand{\arraystretch}{1.3}

% Fill-in field: a thin box for handwritten input
\newcommand{\fillbox}[1][3cm]{\framebox[#1]{\rule{0pt}{2.4ex}\rule{#1}{0pt}}}

\begin{document}

\begin{center}
  {\LARGE\bfseries Code-Review Checklist}\\[0.4em]
  {\large Code-quality check (intra-team review)}
\end{center}
\vspace{0.2em}
\hrule
\vspace{0.6em}

During the intra-team review, each member reads the team's code and scores
it against this checklist. It keeps the review honest and focused on
correctness and clarity, not style.

\textbf{Ground rule:} you are reviewing \emph{correctness and clarity}, not
style taste. ``Not how I would write it'' is not a defect. ``I can't tell what
it does'' is.

\noindent
\textbf{Reviewer:}~\fillbox[4cm] \quad
\textbf{Code under review:}~\fillbox[4cm]

% ============================================================
\section*{A. Can it run?}
% ============================================================
\begin{itemize}[leftmargin=2.4em, itemsep=1pt]
  \item[$\square$] \texttt{python code/run\_all.py} runs end-to-end without errors
  \item[$\square$] README explains how to install deps and run
\end{itemize}

% ============================================================
\section*{B. Is the model right?}
% ============================================================
\begin{itemize}[leftmargin=2.4em, itemsep=1pt]
  \item[$\square$] The RHS $f(t, y)$ matches the report's equations (spot-check 2--3 terms)
  \item[$\square$] Parameters are defined somewhere findable (constants block / config)
  \item[$\square$] Units or scales are noted (e.g.\ stiffness ratio computed)
\end{itemize}

% ============================================================
\section*{C. Are the methods correctly implemented?}
% ============================================================
\begin{itemize}[leftmargin=2.4em, itemsep=1pt]
  \item[$\square$] Explicit Euler: $y + h f(t, y)$ --- one line, not accidentally implicit
  \item[$\square$] RK4: all four stages use the correct arguments ($k_2$ uses $t+h/2$, etc.)
  \item[$\square$] Implicit Euler: Newton solves $w - y_n - h f(t_{n+1}, w) = 0$, not a
    simplified version
  \item[$\square$] Newton Jacobian $J_F = I - h J_f$ (check sign!) matches the
    finite-difference check in the report
  \item[$\square$] No hard-coded step sizes where the code claims adaptivity
  \item[$\square$] Adaptive controller estimates an error and \emph{uses} it to change $h$
\end{itemize}

% ============================================================
\section*{D. Are the numbers believable?}
% ============================================================
\begin{itemize}[leftmargin=2.4em, itemsep=1pt]
  \item[$\square$] There is a comparison against a reference (analytic or high-accuracy)
  \item[$\square$] The observed order is consistent with the claim (slope on log-log)
  \item[$\square$] Mass / energy / invariant conservation is checked where applicable
  \item[$\square$] Any claim like ``stable'' is backed by a stability-region plot or an
    experiment, not just an assertion
\end{itemize}

% ============================================================
\section*{E. Is it maintainable?}
% ============================================================
\begin{itemize}[leftmargin=2.4em, itemsep=1pt]
  \item[$\square$] Functions are small and named by what they do
  \item[$\square$] $f(t, y)$ is swappable (the integrator doesn't know the model)
  \item[$\square$] Magic numbers are named or commented
  \item[$\square$] Comments explain \emph{why}, not just \emph{what}
  \item[$\square$] No dead code (commented-out blocks, unused imports)
\end{itemize}

% ============================================================
\section*{F. Findings}
% ============================================================
\begin{table}[h]
\centering
\small
\begin{tabular}{@{}p{0.8cm}p{3.2cm}p{3.2cm}p{3.2cm}p{3.5cm}@{}}
\toprule
\# & Severity (blocker / major / minor) & Location (file:line) & What's wrong & Suggested fix \\
\midrule
1 & & & & \\
2 & & & & \\
3 & & & & \\
\bottomrule
\end{tabular}
\end{table}

% ============================================================
\section*{G. Top 3 things to fix before Seminar 2}
% ============================================================
1.\ \fillbox[5cm] \quad 2.\ \fillbox[5cm] \quad 3.\ \fillbox[5cm]

% ============================================================
\section*{Reviewer instructions}
% ============================================================
\begin{itemize}[leftmargin=2em, itemsep=2pt]
  \item \textbf{Blocker} = code doesn't run, or produces wrong results (must fix).
  \item \textbf{Major} = a claim in the report is not supported by the code
    (e.g.\ ``adaptive'' but fixed $h$; ``order 4'' but RK4 stages wrong).
  \item \textbf{Minor} = clarity / maintainability (rename, comment, split function).
  \item Fill in findings even if everything passes --- ``no findings'' means you
    did not look carefully.
  \item Give the reviewed team one actionable \emph{positive}: something they
    did well that they should keep doing.
\end{itemize}

\end{document}